\section{Introduction}

Paper B.1 identified where the interstate system is absent: 66.5 million Americans live in counties more than 30 miles from any interstate on-ramp. This paper addresses the complementary problem: where the system is present but overwhelmed. These are not the same places. Coverage gaps are predominantly rural --- the Northern Tier, the Appalachians, the Rural Gulf South. Bottlenecks are predominantly urban --- Atlanta, Chicago, Los Angeles, New Jersey. The investment strategies are different, the political constituencies are different, and the I2.0 interventions are different.

Paper A.1 established the congestion-stress paradox: aggregate scoring based on throughput metrics places short urban connectors (I-110, I-880, I-285) above long transcontinental arteries (I-80, I-95) in the tier rankings, because urban connectors operate at or above capacity while transcontinental routes average congested urban endpoints against free-flowing rural spans. We revisit this paradox through the lens of the American Transportation Research Institute's annual bottleneck rankings \citep{ATRI2024}, which measure actual freight congestion costs at the 100 worst locations in the US highway network.

The question is: do the most economically costly bottlenecks occur on T1 Primary Arteries (strategically important, nationally connected) or on T2 Major Connectors (regionally important, locally congested)? The answer has direct implications for I2.0 investment prioritization. If T1 corridors dominate bottleneck costs, the managed freight lane investment on T1 arteries is the correct priority. If T2 connectors dominate, the immediate investment case is for T2 relief --- improving the connectors that I-75 Atlanta traffic uses to bypass the I-285 interchange.

We find, somewhat paradoxically, that both things are true in different dimensions: T2 connectors account for most bottleneck \textit{locations}, while T1 corridors account for most bottleneck \textit{economic cost}. The resolution is cascade effects: a T1 bottleneck at I-80 near the Bay Area disrupts freight across the entire northern transcontinental route; a T2 bottleneck at I-285 near Atlanta disrupts regional distribution but does not cascade nationally. The same delay on I-80 costs more than the same delay on I-285 because more freight depends on I-80's uninterrupted operation.

\subsection*{Contributions}

\begin{enumerate}
\item Integration of ATRI bottleneck cost data with ROUTE tier classification: first systematic analysis of freight bottleneck costs by strategic tier
\item Quantification of T1 cascade multiplier: T1 bottlenecks impose 1.7× the economic cost per location of T2 bottlenecks
\item Identification of weather bottlenecks (Donner Pass, Snoqualmie Pass) as a distinct bottleneck type not captured in congestion metrics
\item Investment prioritization: managed freight lanes on T1 bottlenecks have higher NPV per dollar than T2 relief, despite lower congestion severity scores
\end{enumerate}
